\documentclass[a4j]{jarticle}
\usepackage{amsmath,amssymb,amsfonts,amscd}
\usepackage[dvipdfmx]{graphicx}
\bibliographystyle{jplain}
\title{演習答案のテンプレート%
    \thanks{本文書を作るにあたって{\tt https://www.overleaf.com/latex/examples/platex-shi-yan/cmhydctsbbfq}をテンプレートとして使わせて頂きました。ありがとうございます。}}
\date{\today}
\author{me}

%% commands
\newcommand{\bvec}[1]{\boldsymbol{#1}}
% 矢印ベクトル対応
%\renewcommand{\bvec}[1]{\vec{#1}}

\begin{document}
\maketitle
\begin{abstract}
答案を書くひとつの方法を説明する。この文書をテンプレートとして{\tt overleaf.com}で作業し、答案 PDFを作成して提出すれば良い。なお、この方法は相応の手間がかかるので、他の作成手段を利用できるならそちらの方が望ましいだろう。
\end{abstract}

\section{方法}
ここで紹介する方法は \TeX を使うこと。これは数式や図を含む文書を作成する有力なものの一つ。いくつか理由を挙げると、数学や物理の文書(論文や教科書)を作成するための十分な機能を有している。 {\tt overleaf.com}のようなウェブ上での作成環境がある。また、できるだけネットワーク接続しないようにしたい場合、パソコン上でも作業できる。ただし、パソコン上での環境準備はそれなりの手間と知識が必要。

\section{例文}
パラグラフに数式を埋めこむなら 
$x$, $y$, $xy$, $x^2y$, $x^{21+3.34}y^{1-0.00001}$.
関数は
平方根 $\sqrt{x^2+y^2}$,
$\cos(x)$, $\sin(x)$, $\arcsin(x)$, $\log 10$, $l\exp(22)$.
分数や微分、偏微分は
$2/3$, $dx/dt$, $\partial f/\partial x$.
添字は $T_{i}$, $S_{ij}$ 等々。
ベクトルは
$\vec{r}$ 等。
時間微分を $\dot{x}$, $\ddot{x}$ と書くこともできる。
$f'(x)$, $f''(x)$, $f^{(2)}(x)$ とかもある。

数式を大きくする例 (いわゆる displaystyle):
ニュートンの運動方程式
$$
m\frac{d^2}{dt^2} \vec{r} =\vec{F}(\vec{r})
$$
分数 $1/3$ を
$$
\frac{1}{3}
$$
と書くことはできる。
微分 $dx/dt$ や偏微分 $\partial f/\partial x$ を
$$
\frac{dx}{dt}
\quad % 間をあける
\frac{\partial f}{\partial x}
$$
と書くこともできる。

図の付けかた。まずお絵描きソフトなりなんなりを使って画像ファイルを生成する必要がある。画像ファイルをアップロードするとこの文書にはめこめる。
\begin{center}
    \includegraphics[width=10cm]{atan}
\end{center}


\section{備考}
作業するのにできるだけネットワーク接続しないようにしたい場合、パソコン上でもできる。ただし、パソコン上での環境準備はそれなりの手間と知識が必要。この場合を考えて、古い処理系でも使えるように、documentclass を jarticle にした。ただし、この場合の図の付けかたは未検証。

\end{document}